\documentclass[11pt]{amsart}
\usepackage{amsmath,amssymb,amsthm,mathtools}
\usepackage[margin=1in]{geometry}
\usepackage{booktabs}
\usepackage[hypertexnames=false]{hyperref}
\hypersetup{hidelinks}
\usepackage[final,expansion=false]{microtype}
\setlength{\parskip}{3pt plus 1pt}
\emergencystretch=2em

\newtheorem{theorem}{Theorem}[section]
\newtheorem{proposition}[theorem]{Proposition}
\newtheorem{lemma}[theorem]{Lemma}
\newtheorem{corollary}[theorem]{Corollary}
\theoremstyle{definition}
\newtheorem{definition}[theorem]{Definition}
\theoremstyle{remark}
\newtheorem{remark}[theorem]{Remark}
\newtheorem{observation}[theorem]{Numerical observation}

\newcommand{\R}{\mathbb R}
\newcommand{\C}{\mathbb C}
\newcommand{\Rel}{\operatorname{Re}}
\newcommand{\Tr}{\operatorname{Tr}}
\newcommand{\supp}{\operatorname{supp}}
\newcommand{\ran}{\operatorname{ran}}
\newcommand{\IT}{I_T}
\newcommand{\PT}{\mathcal P_T}
\newcommand{\Fh}{\widehat F}
\newcommand{\cF}{c_F}
\newcommand{\NT}{\mathcal N_T}
\newcommand{\Aconst}{\mathsf A_T}
\newcommand{\PiT}{\Pi_T}
\newcommand{\lamA}{\lambda_{\mathrm A}}
\newcommand{\Son}{\mathrm{Son}}
\newcommand{\Pih}{\widehat P}

\begin{document}

\title[A semilocal Sonin space and a finite reduction]{A semilocal Sonin space and a finite reduction of window positivity:\\ notes on a construction of Connes and Consani}
\author{David Alejandro Trejo Pizzo}
\address{Independent researcher, Buenos Aires}
\date{October 2026}

\begin{abstract}
Starting from the Sonin space and the function $\varepsilon$ of Connes and Consani, we prove three
things about the positivity of Weil's quadratic form on a window of half-length $T$. First, the
remainder left by the archimedean Sonin trace has infinite negative index on primitive tests as soon as
$T>\frac12\log2$, so the ``one negative direction'' structure of their archimedean theorem stops exactly
at the window $|I|=\log2$. Second, the semilocal Fourier transform $F_S=EF_\infty E^{-1}$ for $S=\{\infty,2\}$ is a
unitary self-adjoint involution whose scattering phase relative to $F_\infty$ reproduces the prime comb
of $2$ exactly; the compressed operator $PF_SP$ is compact but not Hilbert--Schmidt, the semilocal Sonin
projection has an operator gap, and the semilocal analogue of the Connes--Consani trace formula holds with
a remainder density having explicit logarithmic singularities at the powers of $2$, so that the window
form is a positive trace form minus a Hilbert--Schmidt form at every $T$. Third, an archimedean coercivity
theorem with explicit constants on functions satisfying the two Tate conditions blockwise reduces window
positivity to the positivity of an explicit real symmetric matrix of size $2k-2$. No sign is obtained; the results say in which subspace the sign has to be decided.
\end{abstract}

\maketitle

\section{Introduction}

I have been working on this problem for more than ten years, as an amateur, outside academia. Of all the
routes I have tried, the one I kept returning to is yours: the Sonin space, the function $\varepsilon$, the
prolate operators and the semilocal adelic picture that you built with Caterina Consani and Henri
Moscovici. It has a visible structure: one sees where the primes enter and what the archimedean place does on its own.

Let me say where I stand. I think the positivity of the quadratic form on windows of
growing length needs a new idea, and I do not have it. What I do have is a body of
work built directly on your construction: an exact semilocal version of your trace formula, a precise
statement of where your archimedean mechanism stops, a few theorems on the archimedean form that I have
not found in the literature, and a finite reduction that says exactly in which subspace the difficulty
sits. These notes collect that work, with proofs sketched and every statement labelled by what it is: a
theorem, a numerical observation, or an open question. If any of it helps you move the route forward, I
would be glad; if not, I would still like you to have it.

\subsection*{Notation}
Throughout, $T>0$, $\IT=(-T,T)$, and for $F\in L^2(\IT)$ we write $\Fh(\tau)=\int F(t)e^{i\tau t}dt$ and
$\cF(u)=\int F(t)\overline{F(t-u)}dt$ for the autocorrelation, supported in $[-2T,2T]$. Put
\[
m(\tau)=\Rel\psi\big(\tfrac14+\tfrac{i\tau}2\big)-\log\pi,\qquad v_n=\frac{2\Lambda(n)}{\sqrt n},\qquad
\NT=\{n\ge2:\ \log n<2T\},\qquad \Aconst=\sum_{n\in\NT}v_n,
\]
and define the \emph{window form of half-length $T$} and its two pieces,
\begin{equation}\label{eq:form}
\mathfrak a_T(F)=\Gamma(F)+\PiT(F),\qquad
\Gamma(F)=\frac1{2\pi}\int_\R m(\tau)|\Fh(\tau)|^2d\tau,\qquad
\PiT(F)=-\sum_{n\in\NT}v_n\Rel\cF(\log n),
\end{equation}
on the form domain $\{F:\int\log(2+|\tau|)|\Fh|^2<\infty\}$. The two Tate moments are
$A(F)=\int Fe^{t/2}$, $B(F)=\int Fe^{-t/2}$, the \emph{primitive subspace} is $\PT=\{A(F)=B(F)=0\}$, and
$\lamA(T)$ denotes the bottom of $\mathfrak a_T$ on the unit sphere of $\PT$. The complete form
$W_T=\mathfrak a_T+2\Rel(A\bar B)$ is the one the explicit formula writes as a sum over zeros; on $\PT$ the
two agree. In your notation, $g\in C_c^\infty(\R_+^*)$, $f=g*g^*$ and $F(t)=g(e^t)$, so that $f(e^x)=\cF(x)$,
$\widehat g(0)=\int F$, $\widehat g(\mp i/2)$ are the two Tate moments, and $W_\infty(g*g^*)=\Gamma(F)$ by
\cite[(152)--(154)]{CC2021}; the condition $\widehat g(i/2)=0$ of your Theorem~1 is one of the two defining $\PT$; $\PT$ also imposes $\widehat g(-i/2)=0$ and does not impose $\widehat g(0)=0$. The target throughout is $\lamA(T)\ge0$, window by window.
A certificate exists for $T\le0.8$ \cite{Zhu2026}, and $\lamA\ge2.5\cdot10^{-12}$ for
$T\le\frac12\log5=0.8047$ \cite[\S15]{Paper47}, and $\lamA(T)\to0$ at least as fast as
$T^2(k/T)^{2k}$ for every fixed $k$ \cite[Thm.~11.5]{Paper47}. The form $\mathfrak a_T$ is also a L\'evy
energy, $\Gamma(F)=C_0\|F\|^2+\int_0^\infty(\|F\|^2-\Rel\cF(u))K(u)du$ with
$K(u)=2e^{-u/2}/(1-e^{-2u})$ and $C_0=\psi(\frac14)-\log\pi=-5.3722$ \cite[\S2]{Paper47}; this
is the form in which the blocks of \S\ref{sec:coer} are computed.

\subsection*{Summary of results}
Your archimedean theorem \cite[Thm.~1 and Thm.~6.11]{CC2021} writes $W_\infty(g*g^*)$ as the Sonin trace
$\Tr(\vartheta(g)\mathbf S\vartheta(g)^*)\ge0$ plus a remainder which, on the window $|I|=\log2$, has exactly
one negative direction, controlled by $|\widehat g(0)|^2$. We asked what this split does beyond
$\log2$, when primes enter, and whether the semilocal Sonin spaces of \cite{CCquasi,CCM2023,CCM2024} can
absorb them. Briefly:
\begin{itemize}
\item[(i)] the ``one negative direction'' structure stops exactly at $T=\frac12\log2$: for every larger $T$ the
remainder has infinite negative index on $\PT$, whatever bounded archimedean term one subtracts
(Theorem~\ref{thm:nogo});
\item[(ii)] for $S=\{\infty,2\}$ the semilocal Fourier transform is a unitary involution whose scattering phase is
the prime comb of $2$; the semilocal Sonin space absorbs the comb exactly and leaves a Hilbert--Schmidt
remainder at every window, with an explicit singular structure (Theorems~\ref{thm:phase}--\ref{thm:84});
\item[(iii)] the archimedean form is coercive, with explicit constants, on functions that satisfy the two Tate
conditions on each of $k$ consecutive blocks; hence $\lamA(T)>0$ if and only if an explicit real symmetric
matrix of size $2k-2$ is positive definite (Theorems~\ref{thm:B}, \ref{thm:schur}).
\end{itemize}
No sign is obtained. Both pictures point to the same small subspace: in the semilocal picture, a single almost Sonin-orthogonal direction where the window form is a $94\%$ cancellation between two positive quantities (\S\ref{sec:num}); in the archimedean picture, the span of the piecewise Tate exponentials $e^{\pm t/2}\mathbf 1_{B_j}$ (\S\ref{sec:coer}).

\emph{Labels.} Statements called Theorem, Proposition, Lemma or Corollary are proved; the proofs below are sketches; the full arguments are in the author's notes and are available on request. Numerical observations are stated as such; none of them is a certificate.

\section{The Sonin decomposition of the window form}\label{sec:sonin}

Let $\mathbf S$ be the orthogonal projection on the Sonin space $S(1,1)$ of \cite[Def.~4.4]{CC2021} and
$\varepsilon$ the real, even function of \cite[Thm.~4.7, (83)--(84)]{CC2021} (continuous by (84), see \S\ref{sec:semilocal}), so that
$\Tr(\vartheta(f)\mathbf S)=W_\infty(f)+\int f(\rho^{-1})\varepsilon(\rho)d^*\rho$ for every
$f\in C_c^\infty(\R_+^*)$. For $F\in C_c^\infty(\IT)$ put
\[
\Son(F):=\Tr\big(\vartheta(g)\mathbf S\vartheta(g)^*\big)\ \ge0,\qquad
\langle C_\varepsilon F,F\rangle:=\int_{-2T}^{2T}\cF(x)\,\varepsilon(e^{|x|})\,dx ,
\]
where $C_\varepsilon$ is the operator on $L^2(\IT)$ with the continuous kernel $\varepsilon(e^{|x-y|})$; it is
Hilbert--Schmidt.

\begin{theorem}\label{thm:split}
For every $T>0$ and every $F\in C_c^\infty(\IT)$, with no moment condition,
$\mathfrak a_T(F)=\Son(F)-\langle C_\varepsilon F,F\rangle+\PiT(F)$.
\end{theorem}
\begin{proof}
This is (83) read in the additive variable: $W_\infty(f)=\Gamma(F)$ by (152)--(154), $\Tr(\vartheta(f)\mathbf S)=\Son(F)$ by
cyclicity ($\vartheta(g)$ bounded, $\mathbf S\vartheta(g)^*$ Hilbert--Schmidt), and
\[\int f(\rho^{-1})\varepsilon(\rho)\,d^*\rho=\int\cF(x)\,\varepsilon(e^{|x|})\,dx;\]
the prime comb only sees $k\log p<2T$ by the
support of $\cF$.
\end{proof}

For $T\le\frac12\log2$ the comb is empty and the remainder $-C_\varepsilon$ is compact; your \S6 shows, in the
metric of the variable $\xi$ ($F=\frac12\xi-\xi'$), that it has exactly one negative direction at $T=\frac12\log2$:
$\lambda_{\max}(K_I)\approx1.05158$, $\lambda_2\approx0.6865$ (bounded a priori by $0.7722$ in your \S6), with the eigenvector $\zeta$ of the negative direction overlapping the constant by $\langle\zeta|\xi_0\rangle\approx0.94865$; and by Cauchy interlacing the same ``at most one''
holds for every $T<\frac12\log2$. The next statement says that this is the whole range of that mechanism.

\begin{theorem}\label{thm:nogo}
Let $T>\frac12\log2$, let $k\in L^\infty(-2T,2T)$ be real and even, and let $C_k$ be the operator with kernel
$k(x-y)$ on $L^2(\IT)$. For every $N\ge1$ there is an $N$-dimensional subspace of $C_c^\infty(\IT)\cap\PT$ on which
$-\langle C_kF,F\rangle+\PiT(F)<0$ for $F\neq0$. In particular $\mathfrak a_T-\Son$ has infinite negative index on
$\PT$, also on $\PT\cap\{\int F=0\}$, and no inequality
$\mathfrak a_T-\Son\ge-b|v\rangle\langle v|+c(I-|v\rangle\langle v|)$ with $c>0$, nor any with finitely many negative
directions, holds on $\PT$ for any inner product.
\end{theorem}
\begin{proof}
Fix $c_0\in C_c^\infty(-1,1)$ real (with $\int c_0=0$ for the clause on $\int F=0$), $c_w=c_0(\cdot/w)$ and $\beta=c_w''-\frac14c_w$; two integrations by parts give
$\int\beta e^{\pm t/2}=0$, so all translates of $\beta$ are in $\PT$. Put $a=\log2$ and
$h_j=\beta(\cdot-t_j)+\beta(\cdot-t_j-a)$, $j=1,\dots,N$, with the $t_j$ chosen in $(-T+w,T-a-w)$ so that all mutual
differences avoid the finite set $\{\pm\log n\pm\{0,a\}:n\in\NT\}$ by more than $4w$; this is possible because
$2T-a>0$. For $x=\sum c_jh_j$ the only non-vanishing value of $\cF$ at a contact is the inner one of each pair, at
$\log2$, so $\PiT(x)=-v_2\|\beta\|_2^2\sum|c_j|^2=-\sqrt2\log2\,\|\beta\|_2^2\sum|c_j|^2$, while
$|\langle C_kx,x\rangle|\le4N\|k\|_\infty\|\beta\|_1^2\sum|c_j|^2$ and $\|\beta\|_1^2/\|\beta\|_2^2=O(w)$. For small $w$
the form is $\le-\frac12\sqrt2\log2\,\|\beta\|_2^2\sum|c_j|^2$. The last assertion is a dimension count. With
$k=\varepsilon(e^{|\cdot|})$, bounded on $[0,2T]$ since $\sup_{\rho\ge1}|\varepsilon(\rho)|\le\sum_n\lambda(n)^2/(1-\lambda(n)^2)<\infty$ by (84) and Cauchy--Schwarz, this is $\mathfrak a_T-\Son$.
\end{proof}

\begin{remark}
(1) The theorem uses neither the sign of $\mathfrak a_T$ nor the identity of the prime. It does not contradict
$\mathfrak a_T\ge0$: on the test pairs, $\Son$ grows like $w^{-3}\log(1/w)$ and the comb like $w^{-3}$. It says that the
archimedean Sonin trace, subtracted \emph{in full}, is the wrong positive part once a contact is present. (2) Ceding a
fraction restores finiteness: for every $\eta\in(0,1]$, $n_-(\mathfrak a_T-(1-\eta)\Son|_{L^2(\IT)})<\infty$, bounded by the
number of eigenvalues $\ge A'/(A'+L)$ of the time--band operator of half-length $T$ and bandwidth $s_*$, the least $s$
with $\eta m(s)\ge2A'$, $A'=\|C_\varepsilon\|+\Aconst$ and $L=A'-\eta C_0$; since $m(s)=\log(s/2\pi)+O(s^{-2})$ this is of order
$Te^{2A'/\eta}$. The index is finite, but the finite arena it leaves is exponentially large in $1/\eta$.
\end{remark}

\section{The semilocal Sonin space at the prime \texorpdfstring{$2$}{2}}\label{sec:semilocal}

\subsection{The semilocal Fourier transform}
We work in $H=L^2(\R_+,du)$, which is $L^2(\R)^{\mathrm{ev}}$ with your inner product $\frac12\int_\R=\int_0^\infty$
\cite[(16)]{CC2021}; $F_\infty$ is the operator with kernel $2\cos(2\pi ux)$, $P=\mathbf 1_{(0,\lambda]}(u)$ is the cut-off,
$(\vartheta(\rho)\xi)(u)=\rho^{-1/2}\xi(\rho^{-1}u)$ is the scaling action, and $\vartheta(f)=\int f(\rho)\vartheta(\rho)d^*\rho$.
In the logarithmic variable $t=\log u$ (unitarily, $\xi\mapsto e^{t/2}\xi(e^t)$) scalings become translations and we write
$D=-i\,d/dt$. For $h$ on $\R_+^*$, $\vartheta(h)$ is then the convolution by $k(s)=h(e^s)$, which we write $\widehat k(D)$. Let $D_2\xi(u)=\xi(2u)$ and $E=\sum_{k\ge0}D_2^k$, so $E^{-1}=1-D_2$; in Mellin form on $\Rel s=\frac12$,
$E$ is the bounded multiplier $(1-2^{-s})^{-1}$. Following \cite[Prop.~4.1]{CCM2023}, where $F_S\circ\eta_S=\eta_S\circ F_\infty$
with $\eta_S=E$ in this model, the \emph{semilocal Fourier transform} is
\[
F_S:=EF_\infty E^{-1}.
\]
Put $\Pih^S=F_SPF_S$, $Q_1=1-P$, $Q_2=1-\Pih^S$, $A_S:=PF_SP$, and let $\Pi_S$ be the orthogonal projection on
$\ran Q_1\cap\ran Q_2$, the semilocal Sonin space (your $\ker u_{22}$ in \cite{CCquasi}). The level $\lambda$ is a free
parameter; the archimedean formula (83) is the case $\lambda=1$.

\begin{theorem}\label{thm:phase}
In the logarithmic variable, $F_\infty=I\,m_\infty(D)$ with $I$ the reflection $t\mapsto-t$ and $|m_\infty|=1$ (the Tate
factor), and $F_S=F_\infty\,e^{i\varphi(D)}$ with
\[
\varphi'(t)=-\sum_{k\ge1}v_{2^k}\cos(kt\log2),\qquad v_{2^k}=2\log2\cdot2^{-k/2}.
\]
For finite $S$ with $F_S=E_SF_\infty E_S^{-1}$, $E_S=\prod_pE_p$, the phases add: $\varphi_S'=-\sum_{p\in S}\sum_kv_{p^k}\cos(kt\log p)$.
Consequently, for every $\lambda>0$ and every $h$ with $y\mapsto h(e^y)$ Schwartz, $\vartheta(h)(\Pih^S-\Pih_\infty)$, $\Pih_\infty=F_\infty PF_\infty$, is trace class and
\begin{equation}\label{eq:Cprime}
\Tr\,\vartheta(h)\big(\Pih^S-\Pih_\infty\big)=\sum_{p\in S}\sum_{k\ge1}\log p\;p^{-k/2}\big(h(p^k)+h(p^{-k})\big),
\end{equation}
exactly, with no remainder and no dependence on $\lambda$.
\end{theorem}
\begin{proof}
With $q=2^{-1/2}$ and $\mathsf T$ the translation by $\log2$, $E^{-1}=1-q\mathsf T$, $F_\infty EF_\infty=E^*$, so
$F_S=F_\infty U$ with $U=E^*E^{-1}=(1-q\mathsf T)/(1-q\mathsf T^{-1})$, a quotient of conjugates: the symbol is
$e^{i\varphi(t)}$, $\varphi(t)=2\arg(1-qe^{it\log2})=-2\sum_kq^k\sin(kt\log2)/k$, and differentiating gives the comb. This $U$ is
the factor $\rho_2(z)=\gamma_2(z)/\gamma_2(1-z)$ of your quasi-inner function $u=\rho_\infty\prod_p\rho_p$
\cite[Lemma~3.1]{CCquasi}; we only verified the normalization. For \eqref{eq:Cprime} write $U=\sum_ju_j\mathsf T^{j}$ with
$\sum_j(1+|j|)|u_j|<\infty$ (geometric coefficients); then, with $P_a=\mathbf 1_{t<a}$ (the cut-off $P$ at level $\lambda=e^a$), $UP_aU^*-P_a=\sum_{j,l}u_j\bar u_l(P_{a-j\log2}-P_a)\mathsf T^{j-l}$,
each term is trace class against $\widehat k(D)$ by Simon's $\ell^1(L^2)$ criterion \cite[Thm.~4.5]{Simon}, and the trace of the
series is $-\frac1{2\pi}\int\widehat k\,\varphi'$; conjugating by $F_\infty$ reverses $D$. Additivity in $S$ follows from the
identity $Y(\varphi_1+\varphi_2)=e^{i\varphi_1(D)}Y(\varphi_2)e^{-i\varphi_1(D)}+Y(\varphi_1)$ for
$Y(\varphi)=e^{i\varphi(D)}P_ae^{-i\varphi(D)}-P_a$ and invariance of the trace.
\end{proof}

The scattering phase of $F_S$ relative to $F_\infty$ is therefore the prime comb of the window form, coefficient by
coefficient, with $\Lambda$ appearing as a logarithmic derivative; \eqref{eq:Cprime} is the dual form of the classical
phase-shift trace formula and is the local term at $p$ of your semilocal trace formula. It concerns the \emph{linear}
object $1-\Pih^S$; the nonlinear part, the intersection $\ran Q_1\cap\ran Q_2$, is where the rest of this section lives.

\begin{proposition}\label{prop:AS}
\emph{(i)} $F_S$ is a real, unitary, self-adjoint involution and $F_S\vartheta(\rho)F_S=\vartheta(\rho^{-1})$; hence
$A_S=PF_SP$ is compact and self-adjoint, $P\Pih^SP=A_S^2$, and $\|A_S\|<1$ ($\pm1$ are not eigenvalues).
\emph{(ii)} $A_S$ is not Hilbert--Schmidt: $\sum_n\lambda_n(A_S)^2=\infty$. It belongs to the Schatten class
$\mathcal S_p$ for every $p>2$.
\end{proposition}
\begin{proof}
(i) $F_S$ is similar to the involution $F_\infty$ and unitary by Theorem~\ref{thm:phase}; $U$ commutes with the scalings.
Compactness is proved in \cite[Thms.~2.1, 4.4]{CCquasi}; directly, the grouped kernel $\sum_{k\le K}(2c_k-c_{k-1})$,
$c_k=\cos(2\pi2^kux)$, gives Hilbert--Schmidt operators converging to $A_S$ in norm, with
$\|PD_2^kF_\infty E^{-1}P\|\le(1+q)q^k$. If $A_S\xi=\pm\xi$ then $F_S\xi=\pm\xi$ is supported in $(0,\lambda]$, so
$E^{-1}\xi$ and its cosine transform are both compactly supported, forcing $\xi=0$.
(ii) If $A_S$ had an $L^2$ kernel $\kappa$, pairing with $\chi c_m$, $\chi\in C_c^\infty((\delta,\lambda-\delta)^2)$, would
give $\langle\kappa,\chi c_m\rangle\to0$ by Riemann--Lebesgue, while the lacunary structure of the kernel (non-stationary phase
in every term $c_jc_m$, $j\ne m$) gives $\langle\kappa,\chi c_m\rangle=\frac12\iint\chi+O(2^{-m})\ne0$. For $p>2$,
$\|Pc_kP\|_{\mathcal S_p}\le\|Pc_kP\|^{1-2/p}\|Pc_kP\|_{\mathrm{HS}}^{2/p}\le(\tfrac12q^k)^{1-2/p}\lambda^{2/p}$ is summable.
\end{proof}

So the primes do change the spectrum of the band operator: the similarity $A_S=EA_\infty E^{-1}$ is false
(the defect is $PE(1-P)F_\infty PE^{-1}P$, the mass outside the box folded back by the dyadic dilations), and the
superexponential Slepian decay of the archimedean eigenvalues is replaced by a square-divergent sequence. At
$\lambda=\frac12$ the first eigenvalues are $0.988$, $-0.507$ (numerical) against $0.885$, $-0.107$ for $A_\infty$;
at $\lambda=1$, $1-\lambda_1\approx2.6\cdot10^{-6}$ and $\lambda_2\approx-0.9927$ (numerical; stable to the digits shown under basis refinement, higher digits drift),
so the gap in the next theorem is positive but tiny at the level $\lambda=1$.

\subsection{The semilocal trace formula}
Let $\{\xi_n\}$ be a real orthonormal eigenbasis of $A_S$ in $\ran P$, $A_S\xi_n=\lambda_n\xi_n$, and for $\lambda_n\ne0$
put $\zeta_n=(1-P)F_S\xi_n/(1-\lambda_n^2)^{1/2}$ and $\tau_n=\lambda_n(1-\lambda_n^2)^{-1/2}$.

\begin{theorem}\label{thm:gap}
\emph{(i)} In Halmos' decomposition of the pair $(Q_1,Q_2)$ the angles satisfy $\cos^2\theta_n=\lambda_n^2$, the $\zeta_n$
are orthonormal, $Q_1Q_2Q_1=\Pi_S+\sum_n\lambda_n^2|\zeta_n\rangle\langle\zeta_n|$, and
$\operatorname{spec}(Q_1Q_2Q_1|_{\ran Q_1})\subset\{0\}\cup\{\lambda_n^2\}\cup\{1\}$ with $\sup\lambda_n^2=\|A_S\|^2<1$. Hence
$\Pi_S$ is the Riesz projection of $Q_1Q_2Q_1$ at the isolated point $1$, and
$\|(Q_1Q_2Q_1)^n-\Pi_S\|=\|A_S\|^{2n}$.
\emph{(ii)} For $g,f\in C_c^\infty(\R_+^*)$: $\vartheta(g)Q_1Q_2$ is Hilbert--Schmidt; $\vartheta(f)Q_1Q_2Q_1$ and
$\vartheta(f)\Pi_S$ are trace class; and $P\vartheta(f)P\,A_S^2$ is trace class, so that
$\sum_n\lambda_n^2\langle\xi_n|\vartheta(f)\xi_n\rangle$ converges absolutely although $\sum\lambda_n^2=\infty$.
\end{theorem}
\begin{proof}
(i) is the two-projection calculus \cite{Halmos1969} with $P\Pih^SP=A_S^2$ from Proposition~\ref{prop:AS}(i).
(ii) By Theorem~\ref{thm:phase}, after the reflection $I$, $Q_1Q_2=I\,\mathbf 1_{t<-a}M(D)\mathbf 1_{t>a}M(D)^*I$ with
$M=m_\infty e^{i\varphi}$ unimodular, $a=\log\lambda$: the semilocal case is the archimedean one with another
symbol. Then $\widehat k(D)\mathbf 1_{t<-a}M(D)\mathbf 1_{t>a}$ splits into a commutator of Hilbert--Schmidt norm
$\int|k|^2|x|dx$ and a Hankel block of norm $\int|K|^2(-2a-x)_+dx$, $K$ the kernel of $\widehat k(D)M(D)$, finite because $\widehat k\,m_\infty e^{i\varphi}$ is
Schwartz ($m_\infty$ has polynomially bounded derivatives by Stirling, $e^{i\varphi}\in C_b^\infty$). The trace-class
statements follow by $f=g_1*g_2$, your Lemma~D.1 for the commutators, and Dixmier--Malliavin;
$\vartheta(f)\Pi_S=\vartheta(f)Q_1Q_2Q_1\Pi_S$. For the last claim, $P\vartheta(f)(\Pih^S-\Pih_\infty)P=\Theta_{00}(A_S^2-A_\infty^2)+\Theta_{01}X_{10}$, $X_{10}=(1-P)(\Pih^S-\Pih_\infty)P$,
with $\Theta_{01}=P\vartheta(f)(1-P)$ trace class and $A_\infty$ trace class; the left side is trace class by \eqref{eq:Cprime}.
\end{proof}

\begin{theorem}\label{thm:L2}
The series $E_S(f):=\sum_n\tau_n\big(\langle\xi_n|\vartheta(f)\zeta_n\rangle+\langle\zeta_n|\vartheta(f)\xi_n\rangle\big)$ converges
absolutely for every $f\in C_c^\infty(\R_+^*)$ and equals $\int f(\rho^{-1})\varepsilon_S(\rho)d^*\rho$ with
$\varepsilon_S(\rho^{-1})=\varepsilon_S(\rho)$ and, for $\rho\ge1$,
$\varepsilon_S(\rho)=\rho^{1/2}\int_{\lambda/\rho}^\lambda G(\rho x,x)dx$, $G=(1-P)F_SA_S(1-A_S^2)^{-1}P$. For every $R>1$ and
$1\le\rho\le R$,
\[
\varepsilon_S(\rho)=\sum_{1\le j\le\log_2R}\beta_j\,\log^+\frac1{|\rho-2^j|}+O_R(1),\qquad
\beta_j=2^{j/2}\big(1-2^{-j}\big)\frac{\lambda^2}{2\log2},
\]
and $\varepsilon_S(\rho)\sim\frac{\lambda^2}{2\log2}(\rho-1)\log\frac1{\rho-1}$ as $\rho\downarrow1$. Hence
$\varepsilon_S\in L^p_{\mathrm{loc}}(\R_+^*)$ for every $p<\infty$, $\varepsilon_S\notin L^\infty_{\mathrm{loc}}$, and the pointwise
series diverges to $+\infty$ at every $\rho=2^j$, $j\ge1$.
\end{theorem}
\begin{proof}
Absolute convergence: the two cross series are traces of $\Theta_{01}G$ and $\Theta_{10}G^*$ in an orthonormal basis.
Write $G=G_0(1+R)$ with $G_0=(1-P)F_SPF_SP$ and $R=A_S^2(1-A_S^2)^{-1}$, Hilbert--Schmidt by Proposition~\ref{prop:AS}(ii), so
the $R$-part has bounded density. The kernel of $G_0$ is a double lacunary sum of $S(2^ku\mp2^lz)$,
$S(w)=\sin(2\pi\lambda w)/(2\pi w)$, and the lemma $\sum_{k\ge k_0}S(2^kt)=\lambda\,\#\{k\ge k_0:2^k|t|\le1\}+O(1)$, uniform in
$t\neq0$, gives $G_0(\rho x,x)=\frac{\lambda}{2\log2}\sum_i\log^+\frac1{|\rho-2^i|}+O_R(1)$ uniformly in $x\in[\lambda/R,\lambda]$;
integrate over $[\lambda/\rho,\lambda]$. Divergence at $2^j$: the dyadic quasi-covariance of the kernel $K_S$ of $F_S$, $K_S(2^jx,y)=K_S(x,y)-B_j(x,y)$, $B_j$
entire, reduces the pointwise term at $\rho=2^j$ to $\|\mathbf 1_{[2^{-j}\lambda,\lambda]}A_S\|_{\mathrm{HS}}^2$, which is infinite
by the argument of Proposition~\ref{prop:AS}(ii) localized to an interval. The numerical slopes of the three first
singularities, $0.128$, $0.272$, $0.448$ at $\lambda=\frac12$, agree with $\beta_1,\beta_2,\beta_3$ to one percent.
\end{proof}

The non-Hilbert--Schmidt character of $A_S$ thus shows up in exactly one place: the term $m=0$ of the Neumann series of
$(1-A_S^2)^{-1}$, whose density has positive logarithmic peaks at the powers of $2$, of weight growing like $2^{j/2}$. At
$\lambda=1$, $\beta_1=0.510$. There is no atom at $\rho=2$ in $\varepsilon_S-\varepsilon_\infty$: the atom is the comb, which
the next theorem separates.

\begin{theorem}\label{thm:84}
Let $S=\{\infty,2\}$, $f\in C_c^\infty(\R_+^*)$ and $\mathrm{comb}_2(f)=\sum_{k\ge1}\log2\cdot2^{-k/2}\big(f(2^k)+f(2^{-k})\big)$.
\emph{(i)} For every $\lambda>0$,
$\Tr\,\vartheta(f)\Pi_S-E_S(f)=\Tr\,\vartheta(f)\Pi_\infty-E_\infty(f)-\mathrm{comb}_2(f)$.
\emph{(ii)} For $0<\lambda\le1$, the level dictionary
$\Tr\,\vartheta(f)\Pi_\infty^{(\lambda)}=W_\infty(f)+2\log\tfrac1\lambda\,f(1)+\int f(\rho^{-1})\varepsilon_\lambda(\rho)d^*\rho$ holds,
with $\varepsilon_1=\varepsilon$, and therefore
\[
\Tr\,\vartheta(f)\Pi_S=W_\infty(f)-\mathrm{comb}_2(f)+2\log\tfrac1\lambda\,f(1)+\int f(\rho^{-1})\varepsilon_S(\rho)\,d^*\rho .
\]
\end{theorem}
\begin{proof}
(i) For $f$ real with $f(\rho^{-1})=f(\rho)$ one has $[F,\vartheta(f)]=0$ for $F\in\{F_\infty,F_S\}$, and expanding
$\langle\xi_n|F\vartheta(f)F\xi_n\rangle$ with $F\xi_n=\lambda_n\xi_n+(1-\lambda_n^2)^{1/2}\zeta_n$ gives, for each $n$ and without
summing, $\langle\zeta_n|\vartheta(f)\zeta_n\rangle=\langle\xi_n|\vartheta(f)\xi_n\rangle-\tau_n(c_n+c_n')$ with $c_n=\langle\xi_n|\vartheta(f)\zeta_n\rangle$, $c_n'=\langle\zeta_n|\vartheta(f)\xi_n\rangle$, your (87). By Theorem~\ref{thm:gap}(ii) the four series (two diagonal, two cross) converge absolutely, so your passage (85)--(89) can be performed term by term: for each $n$, $\lambda_n^2\langle\xi_n|\vartheta(f)\xi_n\rangle=\lambda_n^2\langle\zeta_n|\vartheta(f)\zeta_n\rangle+\lambda_n^2\tau_n(c_n+c_n')$, and $\lambda_n^2\tau_n+\lambda_n(1-\lambda_n^2)^{1/2}=\tau_n$ collects the two cross coefficients into $E_F(f)$; your (90), $\Tr\vartheta(f)Q_1Q_2Q_1-\Tr\vartheta(f)\Pi_F=\sum_n\lambda_n^2\langle\zeta_n|\vartheta(f)\zeta_n\rangle$, is Theorem~\ref{thm:gap}(i).
Subtracting the two cases, $Q_1Q_2^SQ_1-Q_1Q_2^\infty Q_1=-(1-P)(\Pih^S-\Pih_\infty)(1-P)$, whose trace against $\vartheta(f)$ is
computed blockwise and equals $-\mathrm{comb}_2(f)$ by \eqref{eq:Cprime}. Antisymmetric $f$ give zero on both sides since
$F\Pi_F=\Pi_FF$. (ii) $S(\lambda,\lambda)$ is unitarily $S(1,\lambda^2)$ by a scaling commuting with $\vartheta(f)$, and the
passage from $S(1,b)$ to $S(1,1)$, $b=\lambda^2\le1$, adds $\log\frac1b\,f(1)$ plus an explicit continuous density; combine with
(83) at $\lambda=1$.
\end{proof}

The term $2\log\frac1\lambda\,f(1)=2\log\frac1\lambda\,\|F\|^2$ ($=1.386\|F\|^2$ at $\lambda=\frac12$) is not negligible: without it the semilocal Sonin trace is negative on the near-minimizer of the window form; with it every check we could make passes, including the monotonicity
$\Son_\infty^{(1/2)}-\Son_\infty^{(1)}\succeq0$ for the archimedean forms (numerical margin $\approx1.20$ at $T=\frac12\log3$).

\begin{corollary}\label{cor:split}
Let $\lambda=1$, $S=\{\infty,2\}$. For every $T>0$ and $F\in C_c^\infty(\IT)$,
\[
\Gamma(F)-\sum_{k:\,2^k\in\NT}v_{2^k}\Rel\cF(k\log2)=\Son_{\Pi_S}(F)-\langle C_{\varepsilon_S}F,F\rangle,\qquad
\Son_{\Pi_S}(F):=\Tr\big(\vartheta(g)\Pi_S\vartheta(g)^*\big)\ge0,
\]
and $C_{\varepsilon_S}$, the operator on $L^2(\IT)$ with kernel $\varepsilon_S(e^{|x-y|})$, is Hilbert--Schmidt. In particular, for
$T\le\frac12\log3$, $\mathfrak a_T=\Son_{\Pi_S}-C_{\varepsilon_S}$ on $\PT$: the comb of $2$ is absorbed exactly, and the remainder
is compact, in contrast with Theorem~\ref{thm:nogo}. For $T>\frac12\log3$ the obstruction of Theorem~\ref{thm:nogo} returns
with the contact $3$ (by the proof of that theorem: its test pairs are weakly null and $C_{\varepsilon_S}$ is compact), and the same chain for $S=\{\infty,2,3\}$ absorbs it, with $\varepsilon_S$ now carrying $\log^2$
singularities at the $S$-units $q=a/b$ ($a,b\in\langle2,3\rangle$ coprime), cut off at scale $1/b$ (we expect this form; the lower-order terms are not made explicit).
\end{corollary}

\begin{remark}
From \cite{CCquasi,CCM2023,CCM2024}: $F_S$, the quasi-inner function and the factor $\rho_2$, the compactness of the Hankel
block, the semilocal Sonin space, and the moment and Jacobi structure for $S=\{p,\infty\}$. From \cite{CC2021}: the
archimedean (83)--(90), whose algebra Theorem~\ref{thm:84} transports. What we believe is new: Proposition~\ref{prop:AS}(ii),
Theorem~\ref{thm:gap} as stated (in particular the absolute convergence of the diagonal series), Theorem~\ref{thm:L2} with the
constants $\beta_j$, Theorem~\ref{thm:84} and Corollary~\ref{cor:split}. Nothing here decides a sign: the corollary writes
$\mathfrak a_T$ as a positive trace form minus a compact one, and the singular part of $\varepsilon_S$ is positive as a density (nothing is
asserted about its sign as a form). For growing $S=S(T)=\{\infty\}\cup\{p<e^{2T}\}$ the constants are not uniform: a
heuristic count suggests that $\varepsilon_{S(T)}$ develops singularities of type $1/|\rho-n|$ at every integer $n<e^{2T}$, with
total $L^2$ mass growing exponentially in $T$. This is a conjecture, not a theorem, and it is the point at which the
semilocal remainder stops being ``archimedean plus a small perturbation''.
\end{remark}

\section{Numerical observations on the decomposition}\label{sec:num}

Everything in this section is numerical (Galerkin on $\PT$ with $M\le32$ modes, the metric of $L^2(\IT)$, level $\lambda=1$,
$S=\{\infty,2\}$). The pipeline reproduces your constants: $\varepsilon'(1^+)=22.9965$, and at $T=\frac12\log2$, in the
$\xi$-metric, $\lambda_{\max}(K_I)=1.0504$ against your $1.05158\pm0.0012$, $\lambda_2=0.684$ against your $0.6865$, and
$\langle\xi_0|\zeta\rangle=0.9488$ against $0.94865$.

By Corollary~\ref{cor:split}, $\mathfrak a_T=\Son_{\Pi_S}-C$ on $\PT$ for $T<\frac12\log3$, with $C=\langle C_{\varepsilon_S}\cdot,\cdot\rangle$
and $\Son_{\Pi_S}\succ0$ on $\PT$ (its minimum is $0.0147$ at $T=\frac12\log3$). Put
\[
\rho^*(T):=\sup_{F\in\PT}\frac{C(F)}{\Son_{\Pi_S}(F)},\qquad\text{so that}\qquad \mathfrak a_T\succ0\text{ on }\PT\iff\rho^*(T)<1 .
\]

\begin{observation}\label{obs:rho}
\[
\begin{array}{lcccccc}\toprule
T & \tfrac12\log2 & 0.35 & 0.40 & 0.45 & 0.50 & \tfrac12\log3\\ \midrule
\rho^*(T) & 0.001 & 0.003 & 0.055 & 0.245 & 0.644 & 0.944\\
C(u_T) & 0.003 & 0.006 & 0.021 & 0.026 & 0.023 & 0.017\\
\Son_{\Pi_S}(u_T) & 2.82 & 2.29 & 0.374 & 0.106 & 0.036 & 0.018\\ \bottomrule
\end{array}
\]
Here $u_T$ is the maximizer. The rise of $\rho^*$ from $0$ to $0.94$ comes from the denominator, not from the numerator: $C(u_T)$ stays in
$[0.017,0.026]$ from $T=0.40$ on, while the Sonin mass of the maximizer collapses as $u_T$ aligns with the
direction $s_0(T)$ of least Sonin mass in $\PT$ (overlap $0.998$ at $\frac12\log3$). At $T=\frac12\log3$ the maximizer is the
minimizer $h^*$ of $\mathfrak a_T$ itself: $\Son_{\Pi_S}(h^*)=0.0174$, $C(h^*)=0.0164$, $\mathfrak a_T(h^*)=9.9\cdot10^{-4}$. The
positivity of the window form at its weakest point is a $94\%$ cancellation between two positive quantities of the same size,
not the dominance of the Sonin trace over the remainder. The remainder does have negative directions (the two extreme eigenvalues, $\approx-0.08$ and $\approx-0.04$, are the only
ones stable under refinement), but there $\Son_{\Pi_S}$ wins by a factor $2.5$; and $h^*$ is orthogonal to them.
\end{observation}

\begin{observation}\label{obs:s0}
The direction $s_0(T)$ of least Sonin mass for $\Pi_S$ coincides with the one for $\Pi_\infty$ (overlap $\ge0.989$ on the whole
range) and its minimal mass is the same ($0.0140$ archimedean, $0.0147$ semilocal, at $\frac12\log3$); it is even, with two sign
changes, a ``cosine with negative shoulders'', with $65\%$ of its mass in $|t|<T/2$ and a Fourier peak drifting downward
($11.05$ at $\frac12\log2$, $7.60$ at $\frac12\log3$; a one-mode ``alignment law'' built on the archimedean minimizer, which coincides with $s_0$ to $0.95$, does not reproduce $\lamA(T)$ beyond $T\approx0.45$). It is not your $\zeta$ (overlap $0.03$): $\zeta$ maximizes the
remainder, $s_0$ minimizes the Sonin mass. The prime acts on $s_0$ only through $C$: in the semilocal $s_0$ the balance is $\Delta C=C_S(s_0)-C_\infty(s_0)\approx0.135$ against $\mathfrak a^\infty_T(s_0)\approx0.140$ (the archimedean form in $s_0$), with $C_S(s_0)=+0.010$; in the archimedean minimizer $s_0^\infty$ the same numbers are $C_\infty(s_0^\infty)=-0.140$ and $\mathfrak a^\infty_T(s_0^\infty)=0.153$; over all of $\PT$ the relative loss $\sup(\mathfrak a^\infty_T-\mathfrak a_T)/\mathfrak a^\infty_T$ is
$0.9936$. Concretely: the Fourier peak of $s_0$ sits where $\cos(\tau\log2)>0$ ($\cos(7.6\log2)=0.527$), so
$\Rel c_{s_0}(\log2)=+0.145>0$ and the comb costs $-v_2\cdot0.145=-0.142$, almost all of the archimedean surplus $0.153$ in $s_0^\infty$ ($0.140$ in the semilocal one).
\end{observation}

\section{Archimedean coercivity and a finite reduction}\label{sec:coer}

This part uses only the L\'evy form of $\Gamma$ and the elementary bounds
$\log\frac s{2\pi}-\frac1s-\frac1{2s^2}\le m(s)$ and $m$ increasing on $[0,\infty)$ with $m(0)=C_0$
\cite[Lemma~4.1]{Paper47}. No Sonin space enters.

\begin{definition}
Let $k\ge1$, $h=T/k$, and let $B_1,\dots,B_k$ be the $k$ consecutive intervals of length $2h$ partitioning $\IT$. Put
\[
\mathcal V_k=\Big\{F\in L^2(\IT):\int_{B_j}Fe^{t/2}=\int_{B_j}Fe^{-t/2}=0\ (1\le j\le k)\Big\},\qquad E_k=\PT\ominus\mathcal V_k .
\]
\end{definition}
Then $\mathcal V_k\subset\PT$ is closed, $E_k=\operatorname{span}\{e^{\pm t/2}\mathbf 1_{B_j}\}\cap\PT$ has dimension exactly
$2k-2$, and $\PT=\mathcal V_k\oplus E_k$ orthogonally. Two structural facts about the blocks: the infimum of $\Gamma$ over
functions supported in one block of length $2h\le\log2$ with the two Tate conditions on that block is $\lamA(h)$, the bottom
of a prime-free window of half-length $h$, independently of the position of the block (translation invariance of the
autocorrelation); and for $F,G$ supported on disjoint blocks, $\Gamma(F+G)-\Gamma(F)-\Gamma(G)=-\iint F(x)G(y)K(y-x)$, in which
the rank-one leading term $2e^{-u/2}=2e^{-y/2}e^{x/2}$ of $K$ is killed by either Tate condition, leaving the reduced kernel
$K_1(u)=K(u)-2e^{-u/2}=2e^{-5u/2}/(1-e^{-2u})$, $0<K_1(u)\le u^{-1}$. We have not found this block structure in the literature.

\begin{theorem}\label{thm:A}
For $0<h\le0.1$, $\lamA(h)\ \ge\ \log\frac1h-1.209-h$.
\end{theorem}
\begin{proof}
For $F\in\mathcal P_h$, $\|F\|=1$, the L\'evy form gives $\Gamma(F)=C_0+\int_{2h}^\infty K+\int_0^{2h}(1-\Rel\cF)K$; the tail is
$\log\coth\frac h2+\pi-2\arctan e^h\ge\log\frac2h+\frac\pi2-h$, and $K(u)\ge u^{-1}$ on $(0,2]$. Scaling to
$G(y)=h^{1/2}F(hy)$ on $[-1,1]$, the local integral is $\ge\int_0^2\frac{1-\Rel c_G(v)}vdv=\frac1{2\pi}\int w|\widehat G|^2$ with
$w(s)=\log(2|s|)+\gamma-\operatorname{Ci}(2|s|)$, and $w(s)\ge w(R)\min\{1,s^2/R^2\}$ for $0<R\le\pi$. The two Tate conditions say
$\widehat G(\mp ih/2)=0$, so $\widehat G=(s^2+h^2/4)\widehat\eta$ with $\supp\eta\subset[-1,1]$ and $\eta''\in L^2$; then
$1=\|G\|^2\ge\|\eta''\|^2\ge\frac{\pi^4}4\|\eta\|^2$ (double Poincar\'e inequality on $[-1,1]$), which bounds the low-frequency mass of
$\widehat G$ by $\frac{16(R^2+h^2/4)^3}{27\pi^4R^2}$. With $R=2$, $w(2)=2.1045$ and $C_0+\log2+\frac\pi2=-3.108$, the
constant follows.
\end{proof}
The true value is $\lamA(h)=\log\frac1h-0.510\pm0.001$ on $[0.02,0.1]$ (numerical, Galerkin). The qualitative fact, that
the archimedean form is positive in finite codimension on short windows, is Yoshida's Lemma~3 \cite{Yoshida1992}
(orthogonality to $N$ Fourier modes, no rate) and your Theorem~3.6, Corollary~3.8 and Remark~3.9 in \cite{CC2021} (the archimedean part is
$-2\cdot\mathrm{Id}$ plus Hilbert--Schmidt; positivity for $|I|\le2\cdot0.1404$); the leading term $\log\frac1a$ is observed
numerically for Suzuki's prime-free operator in \cite{Kim2026}. What seems new is the explicit rate with constants under
the two Tate constraints.

\begin{theorem}\label{thm:B}
Let $h=T/k\le1$. For $F\in\mathcal V_k$ and $R>0$, $\|P_RF\|^2\le\delta(Rh)\|F\|^2$ with
$\delta(\theta)=\frac{\cosh^2h}6(\theta^2+\frac{h^2}4)^2$, where $P_R$ is the Fourier projection on $|s|\le R$.
Consequently, for every $F\in\mathcal V_k$ in the form domain,
\[
\Gamma(F)\ \ge\ \Big(\log\frac1h-C_B(h)\Big)\|F\|^2,
\]
with $\kappa=\cosh^2h$, $a=h^2/4$, $\theta_*=(\sqrt{6/\kappa}-a)^{1/2}$ and
$C_B(h)=\log\frac1h-C_0\delta(h)-\frac{2\kappa}3[\Phi_h(\theta_*)-\Phi_h(h)]$, where
$\Phi_h(\theta)=(\frac{\theta^4}4+\frac{a\theta^2}2)\log\frac\theta{2\pi h}-\frac{\theta^4}{16}-\frac{a\theta^2}4-h(\frac{\theta^3}3+a\theta)-\frac{h^2}2(\frac{\theta^2}2+a\log\theta)$
is the primitive of $(\theta^3+a\theta)(\log\frac\theta{2\pi h}-\frac h\theta-\frac{h^2}{2\theta^2})$; $C_B(h)\to\frac14-\log\frac{6^{1/4}}{2\pi}=1.6399$ as $h\to0$, and
$C_B(0.02)=1.657$, $C_B(0.0619)=1.696$, $C_B(0.1)=1.733$.
\end{theorem}
\begin{proof}
On a block with centre $c$, subtract from $G=P_RF$ the element $u$ of $\operatorname{span}\{e^{\pm t/2}\}$ with $u(c)=G(c)$,
$u'(c)=G'(c)$; the remainder $r$ satisfies $r''-\frac r4=G''-\frac G4$ with zero Cauchy data at $c$, whence
$\int_Br^2\le\cosh^2h\frac{h^4}6\|G''-\frac G4\|^2_{L^2(B)}$. Since $F\mathbf 1_B\perp u$, $\langle F,G\rangle=\langle F,r\rangle$;
summing over blocks and using $\|G''-\frac G4\|\le(R^2+\frac14)\|G\|$ gives the leakage bound. Then
$\Gamma(F)/\|F\|^2=\int m\,d\mu=C_0+\int_0^\infty m'(\rho)\mu(|s|\ge\rho)d\rho\ge C_0+\int_0^{\theta_*/h}m'(\rho)(1-\delta(\rho h))d\rho$,
and an integration by parts with the lower bound for $m$ gives the closed form.
\end{proof}
Two moments of Tate per block make $\mathcal V_k$ a high-pass subspace, and $m\sim\log|s|$ does the rest. The theorem
does not separate the block-diagonal part from the coupling: on Tate-orthogonal functions on two adjacent
blocks the bilinear kernel form $\iint F_1F_2K_1(|x-y|)$ has norm \emph{exactly} $\pi$ (Hilbert's inequality for $\le$; a Weyl sequence
$d^{-1/2}\mathbf 1_{[2he^{-M},2h]}(d)$, $d$ the distance to the common endpoint, for $\ge$, the Tate projection costing only
$O(\sqrt M)$), and its contribution to the quadratic form, which counts each unordered pair once, is $\pi/2$. Either way an
operator-norm argument cannot succeed: the essential spectrum of the coupling reaches a constant, while the coercivity gain is
only logarithmic in $k$, and a bound of the form ``diagonal minus coupling'' would need $k\ge70$ at $T=\frac12\log5$ against
$k=25$ from the theorem. In the Galerkin minimizers the coupling costs only $0.17$: the maximizers of the coupling norm are not
the minimizers of the form.

\begin{proposition}\label{prop:prime}
$|\PiT(F)|\le\Aconst\|F\|^2$. More precisely $\|\PiT\|\le\sum_{n\in\NT}v_n\cos\frac\pi{L_n+1}$ with $L_n=\lceil2T/\log n\rceil$,
which is $\Aconst/2$ for $T\le\log2$ ($=0.4901$ for $T\le\frac12\log3$) and $\frac{v_2}{\sqrt2}+\frac{v_3+v_4}2=1.6740$ for
$\log2<T\le\frac12\log5$.
\end{proposition}
\begin{proof}
$\Rel\cF(\ell)=\frac12\langle(S_\ell+S_\ell^*)F,F\rangle$ with $S_\ell$ the shift by $\ell$ compressed to $\IT$; fibred over
$r\in[0,\ell)$, $S_\ell+S_\ell^*$ is the adjacency matrix of a path with at most $L$ vertices, of norm $2\cos\frac\pi{L+1}$.
\end{proof}

\begin{corollary}\label{cor:L1}
$c(T,k):=\inf\{\mathfrak a_T(F):F\in\mathcal V_k,\|F\|=1\}\ \ge\ \log\frac kT-C_B(\frac Tk)-\|\PiT\|$. For instance
$c(\frac12\log3,8)\ge0.486$ ($\dim E_8=14$), $c(\frac12\log5,25)\ge0.094$ ($\dim E_{25}=48$), $c(\frac12\log5,30)\ge0.281$.
For every $T$ there is $k_0(T)\le Te^{\Aconst+C_B}$ with $c(T,k)>0$ for $k\ge k_0(T)$.
\end{corollary}

\begin{theorem}\label{thm:schur}
Fix $T$ and $k$ with $c=c(T,k)>0$. Then $\mathfrak a_T$ is an inner product on $\mathcal V=\mathcal V_k\cap(\text{form domain})$
for which $\mathcal V$ is complete, for every $e\in E_k$ there is a unique $Ke\in\mathcal V$ with
$\mathfrak a_T(Ke,v)=\mathfrak a_T(e,v)$ for all $v\in\mathcal V$, and
\[
\mathfrak a_T(v+e)=\mathfrak a_T(v+Ke)+\mathfrak s_{T,k}(e),\qquad
\mathfrak s_{T,k}(e)=\mathfrak a_T(e)-\sup_{u\in\mathcal V}\big[2\Rel\mathfrak a_T(e,u)-\mathfrak a_T(u)\big].
\]
If $S_{T,k}$ is the matrix of $\mathfrak s_{T,k}$ in an orthonormal basis of $E_k$ and $\lambda_S$ its least eigenvalue, then
$\min\{c,\lambda_S\}/(1+\|K\|)^2\le\lamA(T)\le\min\{c,\lambda_S\}$ when $\lambda_S\ge0$, and $\lamA(T)<0$ when $\lambda_S<0$. Hence $\lamA(T)>0$ if and only if the real symmetric
$(2k-2)\times(2k-2)$ matrix $S_{T,k}$ is positive definite.
\end{theorem}
\begin{proof}
Completeness: on $\mathcal V$, $\mathfrak a_T\ge c\|\cdot\|^2$ and $\mathfrak a_T+(\Aconst-C_0+1)\|\cdot\|^2$ is equivalent to the
weighted norm $\int\log(2+|s|)|\widehat v|^2$. The functional $v\mapsto\mathfrak a_T(e,v)$ is bounded by Cauchy--Schwarz for the
nonnegative form $\mathfrak a_T+(\Aconst-C_0)\langle\cdot,\cdot\rangle$, so $Ke$ exists by Riesz. The identity is the completion of
the square, and the sandwich follows from $\|v\|^2+\|e\|^2\le(1+\|K\|)^2(\|v+Ke\|^2+\|e\|^2)$.
\end{proof}

\begin{remark}
It is an \emph{equivalence}, in an explicit basis of explicit dimension, with the constant $c$ coming from a prime-free
archimedean statement. It is not a gain in precision, and it is not the first finite reduction: sufficient finite reductions in
Legendre bases exist \cite[Thm.~1.1]{Zhu2026}, \cite{Groskin2026}, and the growth $k_0(T)\approx Te^{Ce^T}$ is of the same
doubly exponential order as Zhu's barrier \cite[Thm.~1.4]{Zhu2026}. By the sandwich, $S_{T,k}$ inherits the decay of the margin
in full: at $T=\frac12\log5$, $\lambda_S\approx\lamA\approx3\cdot10^{-12}$ with entries of size $O(1)$ (numerical). Its entries
are not in closed form, since the correction $\mathfrak a_T(Ke)$ is a supremum over $\mathcal V$; Galerkin gives upper bounds for
$S_{T,k}$, and a certified lower bound needs a near-maximizer together with the dual residual divided by $c$, which is
legitimate because $c>0$ is proved. The Galerkin near-minimizers of $\mathfrak a_T$ on $\PT$ have $99\%$ of their mass in $E_k$ for $k\ge9$ ($92\%$ at $T=\frac12\log5$ with $k=5$); this is approximation theory (piecewise Tate exponentials approximate smooth functions to $O(h^2)$), not
arithmetic. The three blocks cut by the orbit of $\log2$ at $T=\frac12\log3$ ($k=3$, $\dim E=4$) do not give a $c$ uniform in $T$: it falls to $0.013$ at
$T=0.69$ (numerical). Related finite-section structure appears in \cite{CvS2025} (simple least eigenvalue with even
eigenfunction) and in Suzuki's screw-function operators on $[-a,a]$ \cite{Suzuki2023,Suzuki2026}; neither uses the Tate
conditions blockwise.
\end{remark}

\section{Open questions}\label{sec:open}

The two halves of this note isolate the same difficulty from two sides without resolving it.
On the semilocal side, the identity $\mathfrak a_T=\Son_{\Pi_S}-C_{\varepsilon_S}$ is exact and its remainder is compact, but at the
weakest direction of the window the two terms cancel to $6\%$, the remainder density has growing positive logarithmic peaks,
and we do not know a lower bound for the Sonin mass $\|\Pi_S\vartheta(g)\|^2_{\mathrm{HS}}$ in terms of $F$ that would see the
comb. On the archimedean side, everything is reduced to the positivity of $S_{T,k}$, whose entries are variational, and whose
least eigenvalue is already $10^{-12}$ at $T=0.805$. Beyond $T\approx0.85$ the least eigenvalue of the window form is at the
floor of double precision in every basis we used; we have no evidence of a sign there in either direction, and we have not tried
to go further than your range and Zhu's.

These are the questions we would most like to ask you.
\begin{enumerate}
\item[(1)] whether the semilocal Sonin space has, at your level $\lambda=1$, a
quantitative version of the operator gap of Theorem~\ref{thm:gap}: the gap $1-\|A_S\|^2$ is $5\cdot10^{-6}$ for $S=\{\infty,2\}$,
and we do not know how it behaves in $S$;
\item[(2)] whether the moment structure of \cite{CCM2024} gives the density
$\varepsilon_S$, or at least its singular coefficients $\beta_j$, in closed form; our derivation goes through the lacunary kernel
and a Neumann series;
\item[(3)] whether a differential operator commuting with $P$ and $\Pih^S$ exists at all: the candidate
$-\mathfrak s^2+\lambda^2N_S$ of \cite{CCM2023} does not commute with $\Pih^S$ for $S=\{\infty,2\}$ (the commutator against
$e^{-5\pi u^2}$ is $-12\pi(30\pi^2u^4-19\pi u^2+1)$ times a nonzero factor), and neither does any $L+\kappa X$ with $L$ a local
differential operator and $X=F_Su^2F_S^{-1}$, because $X$ dilates supports outward across the dyadic scales; so the Halmos
angle operator is, for now, the only operator with that commutation, and it is not differential;
\item[(4)] whether the
direction $s_0$ of Observation~\ref{obs:s0}, a purely archimedean object whose frequency drifts downward with $T$, is something you
recognize from the prolate side.
\end{enumerate}

\begingroup\small
\begin{thebibliography}{99}
\bibitem{CC2021} A.~Connes and C.~Consani, \emph{Weil positivity and trace formula, the archimedean place}, Selecta Math.\ (N.S.) \textbf{27} (2021), no.~4, Paper No.~77; arXiv:2006.13771.
\bibitem{CCquasi} A.~Connes and C.~Consani, \emph{Quasi-inner functions and local factors}, arXiv:2008.10974.
\bibitem{CCM2023} A.~Connes, C.~Consani and H.~Moscovici, \emph{Zeta zeros and prolate wave operators}, arXiv:2310.18423.
\bibitem{CCM2024} A.~Connes, C.~Consani and H.~Moscovici, arXiv:2403.01247 (2024).
\bibitem{CvS2025} A.~Connes and W.~D.~van Suijlekom, \emph{Quadratic forms, real zeros and echoes of the spectral action}, arXiv:2511.23257.
\bibitem{Groskin2026} A.~Groskin, \emph{A finite Guinand--Weil dictionary and archimedean tail order for the truncated Weil quadratic form}, arXiv:2607.02828.
\bibitem{Halmos1969} P.~R.~Halmos, \emph{Two subspaces}, Trans.\ Amer.\ Math.\ Soc.\ \textbf{144} (1969), 381--389.
\bibitem{Kim2026} Kim et al., arXiv:2607.24830 (2026).
\bibitem{Simon} B.~Simon, \emph{Trace Ideals and Their Applications}, 2nd ed., Amer.\ Math.\ Soc., 2005.
\bibitem{Suzuki2023} M.~Suzuki, \emph{Aspects of the screw function corresponding to the Riemann zeta function}, J.~Lond.\ Math.\ Soc.\ (2) \textbf{107} (2023), 1846--1888; arXiv:2206.03682.
\bibitem{Suzuki2026} M.~Suzuki, \emph{Weil's quadratic form via the screw function}, arXiv:2606.09096.
\bibitem{Paper47} D.~A.~Trejo Pizzo, \emph{A dual theory of Weil positivity at finite support: outside-window certificates, non-attainment, and the explicit formula as the dual variable}, preprint, 2026 (available from the author).
\bibitem{Yoshida1992} H.~Yoshida, \emph{On Hermitian forms attached to zeta functions}, in: \emph{Zeta Functions in Geometry}, Adv.\ Stud.\ Pure Math.\ \textbf{21}, Kinokuniya, Tokyo, 1992, 281--325.
\bibitem{Zhu2026} X.~Zhu, \emph{Weil positivity in compact windows: a finite reduction, certified two-sided bounds, and a Landau--Widom decay law}, arXiv:2608.24827.
\end{thebibliography}
\endgroup

\end{document}
